%------------------------------------------------
%   PACKAGES
%------------------------------------------------
\usepackage{xcolor}
\usepackage{tikz}
\usetikzlibrary{shadings}
\usepackage[most]{tcolorbox}

% \EnableBpAbbreviations

\usefonttheme{professionalfonts}


%------------------------------------------------
%   COLORS & TEXT
%------------------------------------------------
\definecolor{myviolet}{RGB}{240, 121, 242}
% \definecolor{lightviolet}{RGB}{255, 200, 255}
\definecolor{lightviolet}{RGB}{255, 179, 255}
\definecolor{lighterviolet}{RGB}{252, 220, 252}


\setbeamercolor{normal text}{fg=black,bg=myviolet}

\setbeamertemplate{navigation symbols}{}

%------------------------------------------------
%   GLAVNI NASLOV
%------------------------------------------------

% Rainbow-bordered title box (reuses rainbow math border)
\tcbset{
    rainbow title/.style={
        enhanced,
        colback=lightviolet,     % background inside the box
        coltext=black,
        arc=6pt,
        outer arc=6pt,
        boxrule=2pt,
        opacityback=1,
        frame code={
            % draw the rainbow border (same as rainbow math)
            \shade[shading=rainbowshading, rounded corners=6pt]
                (frame.south west) rectangle (frame.north east);
        },
        interior code={
            % draw inner light-violet box
            \fill[lightviolet, rounded corners=4pt]
                (interior.south west) rectangle (interior.north east);
        },
        left=12pt, right=12pt, top=10pt, bottom=10pt,
    }
}

% Title page with rainbow-bordered box
\defbeamertemplate*{title page}{rainbowstyle}{
    \begin{center}
        \begin{tcolorbox}[rainbow title]
            \centering
            {\ttfamily\bfseries\Huge\inserttitle}
        \end{tcolorbox}

        \vspace{1em}

        {\large\insertauthor}

        \vspace{0.5em}

        {\small\insertdate}
    \end{center}
}



% \setbeamercolor{frametitle}{bg=lightviolet, fg=black}
% \setbeamertemplate{frametitle}{
%     \begin{beamercolorbox}[ht=1.8em,wd=\paperwidth]{frametitle}
%         \centering
%         {\ttfamily\bfseries\Large\insertframetitle}
%     \end{beamercolorbox}
% }
%------------------------------------------------
%   FRAMETITLE
%------------------------------------------------

\usepackage{tikz}
\usepackage[most]{tcolorbox}

\makeatletter
\tcbset{
    rainbow frametitle/.style={
        % enhanced,
        colback=lightviolet,
        coltext=black,
        % outer arc=0pt,
        boxrule=0pt,
        opacityback=1,
        left=0pt, right=0pt, top=10pt, bottom=10pt,
        width=\paperwidth-8pt,
        frame code={
        },
        overlay={
            % position the box absolutely at the top of the page
            \path[shift={(current page.north west)}] 
                (0,0) rectangle (\paperwidth,-\ht\strutbox);
        }        
    }
}

\setbeamertemplate{frametitle}{
    % use a tcolorbox absolutely positioned
    \begin{tikzpicture}[remember picture,overlay]
        \node[anchor=north west, xshift=-1pt, yshift=1pt] at (current page.north west) {
            \begin{tcolorbox}[rainbow frametitle]
                \centering
                {\ttfamily\bfseries\Large\insertframetitle}
            \end{tcolorbox}
        };
    \end{tikzpicture}
}
\makeatother





% \setbeamertemplate{frametitle}{width={\dimexpr\paperwidth-2\slideinset\relax}}


% \setbeamertemplate{frametitle}{%
%     \noindent\hspace*{\slideinset}% place at inner left edge (do not cover outer border)
%     \begin{tcolorbox}[rainbow frametitle,width={\dimexpr\paperwidth-2\slideinset\relax}]
%         	tfamily\bfseries\large\insertframetitle
%     \end{tcolorbox}
% }





%------------------------------------------------
%   RAINBOW SHADING
%------------------------------------------------
\pgfdeclarehorizontalshading{rainbowshading}{100bp}{
	color(0bp)=(red);
    color(20bp)=(orange);
    color(40bp)=(yellow);
    color(60bp)=(green);
    color(80bp)=(blue);
    color(100bp)=(violet)
}

%------------------------------------------------
%   SLIDE BACKGROUND (RAINBOW BORDER)
%------------------------------------------------
\usebackgroundtemplate{
\begin{tikzpicture}[remember picture,overlay]

    \fill[myviolet]
        (current page.south west) rectangle (current page.north east);

    \shade[shading=rainbowshading]
        (current page.south west) rectangle (current page.north east);

    \fill[myviolet]
        ([xshift=4pt,yshift=4pt]current page.south west)
        rectangle
        ([xshift=-4pt,yshift=-4pt]current page.north east);

\end{tikzpicture}
}



%------------------------------------------------
%   RAINBOW-BORDERED MATH ENVIRONMENT
%------------------------------------------------
\tcbset{
    rainbow math/.style={
        enhanced,
        colback=white,
        % colback=lightviolet,
        coltext=black,
        arc=6pt,                            % rounded corners
        outer arc=6pt,
        boxrule=2pt,
        opacityback=1,
        frame code={
            % draw the rainbow border
            \shade[shading=rainbowshading, rounded corners=6pt]
                (frame.south west) rectangle (frame.north east);
        },
        interior code={
            % draw inner white box
            \fill[lighterviolet, rounded corners=4pt]
                (interior.south west) rectangle (interior.north east);
        },
        left=10pt, right=10pt, top=8pt, bottom=8pt,
    }
}

\newenvironment{whitemath}
{\begin{tcolorbox}[rainbow math]}
{\end{tcolorbox}}


